\documentclass[12pt]{extarticle}

\usepackage{tabularx}
\usepackage{amsmath}
\usepackage{extsizes}
\usepackage{graphicx}
\usepackage{amssymb}
\usepackage[english,russian]{babel}
\usepackage[utf8]{inputenc}
\usepackage[a4paper]{geometry}
\geometry{top=20mm,bottom=30mm,left=20mm,right=15mm}
\usepackage[T1,T2A]{fontenc}
\usepackage{setspace}
\usepackage{titlesec}
\usepackage{indentfirst}
\usepackage{caption}
\renewcommand{\baselinestretch}{1.5}
\usepackage{spverbatim}
\usepackage{listings}
\usepackage{xcolor}
\usepackage{tabularray}
\lstset{extendedchars=\true}

\definecolor{codebg}{gray}{0.96}
\definecolor{comment}{gray}{0.45}
\definecolor{keyword}{RGB}{0, 75, 150}
\definecolor{string}{RGB}{150, 80, 30}
\definecolor{number}{gray}{0.5}
\definecolor{identifier}{rgb}{0,0,0}

\lstdefinestyle{coolprint}{
  basicstyle   = \ttfamily\footnotesize,
  keywordstyle = \bfseries\color{keyword},
  commentstyle = \itshape\color{comment},
  stringstyle  = \color{string},
  identifierstyle = \color{identifier},
  backgroundcolor = \color{codebg},
  numbers      = left,
  numberstyle  = \footnotesize\color{number},
  numbersep    = 5pt,
  frame        = single,
  framerule    = 0.8pt,
  framesep     = 3pt,
  breaklines   = true,
  breakatwhitespace = true,
  tabsize      = 4,
  showspaces   = false,
  showstringspaces = false,
  showtabs     = false,
  xleftmargin  = 10pt,
  xrightmargin = 5pt,
  aboveskip    = \medskipamount,
  belowskip    = \medskipamount,
  captionpos   = b,
}
\lstset{style=coolprint}

\usepackage{ragged2e}
\justifying

\titleformat{\section}{\centering\bfseries}{\thesection}{1em}{}

\begin{document}

\begin{center}
  \textbf{МИНОБРНАУКИ РОССИИ} \\
  \textbf{Санкт-Петербургский государственный} \\
  \textbf{электротехнический университет} \\
  \textbf{«ЛЭТИ» им. В.И. Ульянова (Ленина)} \\
  \textbf{Кафедра САПР} \\
\end{center}

\vspace{0.25\textheight}

\begin{center}
\textbf{Практическая работа №1} \\
\textbf{по дисциплине} \\
\textbf{«Объектно-ориентированное программирование»} \\
\end{center}

\vspace*{\fill}

\begin{center}
\begin{tabularx}{\textwidth}{ l X r }
 Студент гр. 5352 & \rule{8cm}{0.2mm} & Cоболевский М.А. \\
 Преподаватель & \rule{8cm}{0.2mm} & Кулагин М.В. \\
\end{tabularx}
\end{center}

\vspace{3em}

\begin{center}
Санкт-Петербург\\
\the\year
\end{center}

\thispagestyle{empty}

\pagebreak

\section{Задание}

Разработать консольное приложение, которое моделирует работу небольшой службы доставки.

Программа работает только в памяти: внешние файлы, база данных, сетевые запросы и графический интерфейс не являются обязательными.

Пользователь должен иметь возможность создать заказ, выбрать способ доставки, назначить курьера, просмотреть список заказов и переводить заказ по допустимым статусам. В программе должны быть предусмотрены корректная обработка ошибочного ввода и понятные сообщения для пользователя.

\section{Спецификация программы}

\includegraphics[width=\textwidth]{assets/delivery-service-uml}

\section{Описание интерфейса пользователя программы}

\section{Текст программы}

\renewcommand{\baselinestretch}{1}
\input{assets/source}
\renewcommand{\baselinestretch}{1.5}

\section{Выводы}

\end{document}
